\section{Doing useful work with several possible installations}\label{sec:work}
A safe release plan in every possible installation need not yield a safe release plan \emph{for the team}. The first assertion chooses a plan after learning the world. The second needs one policy which can branch only on reports it will actually receive. Permanent refusal can satisfy a safety rule while failing the assigned job. We therefore specify which service must be delivered on a positive verdict and what a negative verdict authorizes.

For a finite class, Theorem~\ref{ext:t29j} tests verdict correctness and actual service together. At every reached node the surviving state set must be nonempty and physically safe. At a positive leaf all its worlds must have the specified target and the required useful goal must already have been achieved in every surviving state within the resource and time limits. At a negative leaf all survivors must have the negative target. For an action chosen before the leaf, every supported successor must have a successful continuation, and the action must be a common callable literal. A finite induction from terminal leaves constructs the policy, or returns an obstruction within that interface. This extends justified choice in Theorems~\ref{thm:t28}--\ref{thm:t29}.

\paragraph{A receipt can pay for itself.}
Suppose an inventory job returns the actual route token as an ordinary work product. Completing it can enable the next required job on a common budget without a purchased audit. Hide the token while preserving the actual authority, and the two later routes can require different first actions at the very same controller record. Knowing that both installations are individually serviceable does not repair this loss.

The correlated serial construction makes this argument computable (Theorem~\ref{ext:t30cg}). An installation is a path through a layered graph, so choices in different job stages can be linked. The target compares the minimum work in the installed path to a deadline. Each mandatory slow job adds its declared work surcharge to the common fast-job base. Cap the accumulated surcharge at one unit above the remaining allowance, where further excess makes no difference to the verdict. For each surviving current edge keep the smallest and largest such prefix scores actually attained. A local read filters edges. A genuine work receipt advances along compatible neighboring edges and updates both extrema with the verified local increment. Monotonicity of the capped addition and attainment preserve the extrema needed for the threshold at every step. The recursion also retains the observed meter, phase, remaining work and decision horizon. It minimizes the worst further price among lawful reads and work moves with successful continuations. The calculation is pseudo-polynomial in the stated numerical bounds at fixed explicit edge width. Unrelated hidden memory, overlapping history effects or a menu that reveals extra information require a fresh sufficient-state argument. Separate static and product results are Theorems~\ref{ext:t30c} and~\ref{ext:t30s}.

This distinction matters for software work. Building an artifact can produce a genuine receipt for the endpoint that will receive it. A simulated receipt may predict the answer while lacking the authority or work effect of the actual build. The experiment in Section~\ref{sec:experiments} includes a case where suppressing one real work observation eliminates the common service policy.
